\SourcePage{551}{554}
\section{Motion in a uniform magnetic field}

For a particle in a constant uniform magnetic field, we shall determine the energy levels (L.~D.~Landau, 1930).


For a uniform field, it is convenient here to choose the vector potential not in the form (111.7), but as follows:
\begin{equation}
 A_x=-Hy,\qquad A_y=A_z=0.
 \tag{112.1}
\end{equation}
(the $z$ axis is chosen along the field). The Hamiltonian then takes the form
\begin{equation}
 \hat H=\frac1{2m}\left(\hat p_x+\frac{eH}{c}y\right)^2+\frac{\hat p_y^2}{2m}+\frac{\hat p_z^2}{2m}-\frac\mu s\hat s_zH.
 \tag{112.2}
\end{equation}

First observe that the operator $\hat s_z$ commutes with the Hamiltonian (since the latter contains no operators for the other spin components). Thus the $z$ projection of the spin is conserved, and $\hat s_z$ may therefore be replaced by its eigenvalue $s_z=\sigma$. The dependence of the wave function on spin is then immaterial, and $\psi$ in the Schrödinger equation may be understood as an ordinary function of the coordinates. For this function we have
\begin{equation}
 \frac1{2m}\left[\left(\hat p_x+\frac{eH}{c}y\right)^2+\hat p_y^2+\hat p_z^2\right]\psi-\frac\mu s\sigma H\psi=E\psi.
 \tag{112.3}
\end{equation}

The coordinates $x$ and $z$ do not occur explicitly in the Hamiltonian of this equation. Hence the operators $\hat p_x$ and $\hat p_z$ (differentiation with respect to $x$ and $z$) also commute with it; in other words, the $x$- and $z$-components of the generalized momentum are conserved. We therefore seek $\psi$ in the form


\begin{equation}
 \psi=\exp\left[\frac i\hbar(p_xx+p_zz)\right]\chi(y).
 \tag{112.4}
\end{equation}

The eigenvalues $p_x$ and $p_z$ range over all values from $-\infty$ to $\infty$. Since $A_z=0$, the $z$ component of the generalized momentum is the same as the corresponding component of the ordinary momentum $mv_z$. The particle's velocity along the field can thus have any value; one may say that motion parallel to the field is ``not quantized.''

Substitution of (112.4) into (112.3) gives the following equation for $\chi(y)$:
\begin{equation}
 \chi''+\frac{2m}{\hbar^2}\left[\left(E+\frac{\mu\sigma}{s}H-\frac{p_z^2}{2m}\right)-\frac m2\omega_H^2(y-y_0)^2\right]\chi=0,
 \tag{112.5}
\end{equation}

\SourcePage{552}{555}
where $y_0=-cp_x/eH$ and
\begin{equation}
 \omega_H=\frac{|e|H}{mc}
 \tag{112.6}
\end{equation}
have been introduced.

Equation (112.5) has the same form as the Schrödinger equation (23.6) for a linear oscillator of frequency $\omega_H$. It follows at once that the quantity in parentheses in (112.5), which plays the role of the oscillator energy, can take the values $(n+1/2)\hbar\omega_H$, where $n=0,1,2,\ldots$

We thus obtain the following energy levels for a particle in a uniform magnetic field:


\begin{equation}
 E=\left(n+\frac12\right)\hbar\omega_H+\frac{p_z^2}{2m}-\frac{\mu\sigma}{s}H.
 \tag{112.7}
\end{equation}
The first term in this expression supplies the discrete energies associated with motion in the plane normal to the field; these are called the \emph{Landau levels}. For an electron, $\mu/s=-|e|\hbar/mc$, and (112.7) becomes


\begin{equation}
 E=\left(n+\frac12+\sigma\right)\hbar\omega_H+\frac{p_z^2}{2m}.
 \tag{112.8}
\end{equation}

The eigenfunctions $\chi_n(y)$ belonging to the energy levels (112.7) follow from (23.12) after the corresponding change of notation ($a_H=\sqrt{\hbar/m\omega_H}$):


\begin{equation}
 \chi_n(y)=\frac1{\pi^{1/4}a_H^{1/2}\sqrt{2^nn!}}\exp\left(-\frac{(y-y_0)^2}{2a_H^2}\right)H_n\left(\frac{y-y_0}{a_H}\right).
 \tag{112.9}
\end{equation}

In classical mechanics, a particle moves in a circle with a fixed center in the plane perpendicular to the field $\mathbf H$ (the $xy$ plane). The quantity $y_0$, conserved in the quantum problem, corresponds to the classical $y$ coordinate of the circle's center. The quantity $x_0=(cp_y/eH)+x$ is conserved as well (its operator is readily verified to commute with the Hamiltonian (112.2)). It corresponds to the classical $x$ coordinate of the circle's center\footnote{Indeed, for classical motion on a circle of radius $cmv_t/eH$ ($v_t$ is the projection of the velocity onto the $xy$ plane; see Vol.~II, \S~21), we have
\[
 y_0=-cp_x/eH=-(cm/eH)v_x+y.
\]
It is evident from this expression that $y$ is the coordinate of the circle's center. The other coordinate is
\[
 x_0=(cm/eH)v_y+x=cp_y/eH+x.
\]}. 

\SourcePage{553}{556}
The operators $\hat x_0$ and $\hat y_0$, however, do not commute with one another, so the coordinates $x_0$ and $y_0$ cannot simultaneously have definite values.

Since (112.7) does not involve $p_x$, whose values form a continuum, every energy level has a continuous degeneracy. This degeneracy becomes finite, however, when the motion in the $xy$ plane is confined to a large but finite area $S=L_xL_y$. The number of distinct (now discrete) values of $p_x$ in an interval $\Delta p_x$ is $(L_x/2\pi\hbar)\Delta p_x$. All $p_x$ for which the orbit centre lies within $S$ are allowed (the orbit radius is neglected in comparison with the large length $L_y$). The conditions $0<y_0<L_y$ give $\Delta p_x=eHL_y/c$. Consequently, for fixed $n$ and $p_z$, the number of states is $eHS/2\pi\hbar c$. If the region of motion is also bounded along the $z$ axis, with length $L_z$, then an interval $\Delta p_z$ contains $(L_z/2\pi\hbar)\Delta p_z$ possible values of $p_z$, and the number of states in that interval is


\begin{equation}
 \frac{eHS}{2\pi\hbar c}\frac{L_z}{2\pi\hbar}\Delta p_z=\frac{eHV}{4\pi^2\hbar^2c}\Delta p_z.
 \tag{112.10}
\end{equation}

For an electron there is an additional degeneracy: the energy levels (112.8) are the same for states with quantum numbers $n$, $\sigma=1/2$ and $n+1$, $\sigma=-1/2$.



\subsection*{Problems}

\ProblemHeading\textbf{1.} Construct the electron wave functions for states in a uniform magnetic field with definite values of momentum and angular momentum along the field direction.


\SolutionHeading In cylindrical coordinates $\rho$, $\varphi$, $z$, with the $z$ axis along the field, the vector potential in the gauge (111.7) has components $A_\varphi=H\rho/2$, $A_z=A_\rho=0$, and the Schrödinger equation is\footnote{We write the electron charge as $e=-|e|$ and denote its mass here by $M$, to distinguish it from the angular momentum $m$. The term containing the particle's spin, which is inessential for this problem, is omitted.}
\begin{equation}
 -\frac{\hbar^2}{2M}\left[\frac1\rho\frac\partial{\partial\rho}\left(\rho\frac{\partial\psi}{\partial\rho}\right)+\frac{\partial^2\psi}{\partial z^2}+\frac1{\rho^2}\frac{\partial^2\psi}{\partial\varphi^2}\right]-\frac{i\hbar\omega_H}{2}\frac{\partial\psi}{\partial\varphi}+\frac{M\omega_H^2}{8}\rho^2\psi=E\psi.
 \tag{1}
\end{equation}
We seek a solution in the form
\[
 \psi=\frac{e^{im\varphi}}{\sqrt{2\pi}}e^{ip_zz/\hbar}R(\rho)
\]
and obtain for the radial function the equation
\[
 \frac{\hbar^2}{2M}\left(R''+\frac1\rho R'-\frac{m^2}{\rho^2}R\right)+\left[E-\frac{p_z^2}{2M}-\frac{M\omega_H^2}{8}\rho^2-\frac{\hbar\omega_Hm}{2}\right]R=0.
\]

\SourcePage{554}{557}
Introducing the new independent variable $\xi=(M\omega_H/2\hbar)\rho^2$, we rewrite this equation as
\[
 \xi R''+R'+\left(-\frac\xi4+\beta-\frac{m^2}{4\xi}\right)R=0,\qquad
 \beta=\frac1{\hbar\omega_H}\left(E-\frac{p_z^2}{2M}\right)-\frac m2.
\]
As $\xi\to\infty$, the required function behaves as $e^{-\xi/2}$, while as $\xi\to0$ it behaves as $\xi^{|m|/2}$. We therefore seek a solution of the form
\[
 R(\xi)=e^{-\xi/2}\xi^{|m|/2}w(\xi)
\]
and obtain for $w(\xi)$ the confluent hypergeometric function equation:
\[
 w=F\left\{-\left(\beta-\frac{|m|+1}{2}\right),|m|+1,\xi\right\}.
\]
For the wave function to be finite everywhere, $\beta-(|m|+1)/2$ must be a nonnegative integer $n_\rho$. The energy levels are then
\[
 E=\hbar\omega_H\left(n_\rho+\frac{|m|+m+1}{2}\right)+\frac{p_z^2}{2M},
\]
which is equivalent to (112.7). The radial wave functions belonging to these levels are
\begin{equation}
 R_{n_\rho m}(\rho)=\frac1{a_H^{1+|m|}}\left[\frac{(|m|+n_\rho)!}{2^{|m|}n_\rho!|m|!^2}\right]^{1/2}\exp\left(-\frac{\rho^2}{4a_H^2}\right)\rho^{|m|}F\left(-n_\rho,|m|+1,\frac{\rho^2}{2a_H^2}\right),
 \tag{2}
\end{equation}
where $a_H=\sqrt{\hbar/M\omega_H}$; the functions are normalized by $\int_0^\infty R^2\rho\,d\rho=1$. Here the hypergeometric function is a generalized Laguerre polynomial.

\ProblemHeading\textbf{2.} Determine the lowest energy level associated with a bound state of an electron in a shallow potential well $U(r)$ ($|U|\ll\hbar^2/ma^2$, where $a$ is the range of the forces in the well) when a uniform magnetic field is also applied (Yu.~A.~Bychkov, 1960).

\SolutionHeading The stated condition on the field $U(r)$ ensures that, in the absence of a magnetic field, perturbation theory is applicable to it; the well then has no bound states (see \S~45). When a magnetic field is present as well, $U(r)$ can be treated as a perturbation only for motion in the plane transverse to $\mathbf H$, whose (discrete) type of energy spectrum is not changed by adding $U$. The nature of the motion along $\mathbf H$, however, changes: as will be seen, unbounded motion becomes bound, and the continuous spectrum becomes discrete. The well's field therefore cannot be treated by perturbation theory for this motion.

Accordingly, on separating variables in the Schrödinger equation (equation (1) of the preceding problem, supplemented by the term $U\psi$ on the left), we retain the radial functions $R(\rho)$ in the previous form (2); the lowest level corresponds to the quantum numbers $n_\rho=m=0$. Substituting $\psi=R_{00}(\rho)\chi(z)$ into the Schrödinger equation, then multiplying it by $R_{00}(\rho)$ and integrating over $\rho\,d\rho$, we obtain for $\chi(z)$
\begin{equation}
 -\frac{\hbar^2}{2m}\chi''+\overline U(z)\chi=\varepsilon\chi,
 \tag{3}
\end{equation}

\SourcePage{555}{558}
where $\varepsilon=E-\hbar\omega_H/2$ and
\[
 \overline U(z)=\int_0^\infty U\left(\sqrt{z^2+\rho^2}\right)R_{00}^2(\rho)\rho\,d\rho
\]
($m$ once again denotes the particle's mass). In form, this equation is the Schrödinger equation for one-dimensional motion in the potential well $\overline U(z)$, with $\varepsilon$ as the energy of that motion. We may therefore use the result of Problem~1 in \S~45, according to which the discrete energy level is
\begin{equation}
 \varepsilon=-\frac m{2\hbar^2}\left[\int_{-\infty}^{\infty}\overline U(z)\,dz\right]^2=-\frac m{2\hbar^2}\left[\int_{-\infty}^{\infty}\int_0^\infty U\left(\sqrt{z^2+\rho^2}\right)R_{00}^2(\rho)\rho\,d\rho\,dz\right]^2.
 \tag{4}
\end{equation}

The wave function $R_{00}(\rho)$ decays over distances $\rho\sim a_H$. If the magnetic field is so weak that $a_H\gg a$, the integral over $d\rho$ is determined by the region $\rho\lesssim a$, in which one may set $R_{00}(\rho)\approx R_{00}(0)=1/a_H$. Then
\begin{equation}
 \varepsilon=-\frac{me^2H^2}{8\pi^2\hbar^4c^2}\left(\int U(r)\,dV\right)^2
 \tag{5}
\end{equation}
($dV=2\pi\rho\,d\rho\,dz\to4\pi r^2\,dr$). In the opposite case of a strong magnetic field, where $a_H\ll a$, the integral in (4) is determined by the region $\rho\lesssim a_H$, in which one may put $U(\sqrt{\rho^2+z^2})\approx U(z)$. The integration over $d\rho$ then reduces to the normalization integral for $R_{00}$ and equals 1, so that
\begin{equation}
 \varepsilon=-\frac{2m}{\hbar^2}\left(\int_0^\infty U(z)\,dz\right)^2.
 \tag{6}
\end{equation}
In both cases, an estimate of the integrals shows that $\varepsilon\ll\hbar\omega_H$.

\ProblemHeading\textbf{3.} Determine the energy levels of a hydrogen atom in a magnetic field so strong that $a_H\ll a_B$, where $a_B$ is the Bohr radius (R.~J.~Elliott, R.~Loudon, 1960).

\SolutionHeading Under the stated condition, $\hbar\omega_H\gg me^4/\hbar^2$, the effect of the nuclear Coulomb field on the electron's motion in the plane transverse to $\mathbf H$ may be treated as a small perturbation. We therefore return to the situation considered in Problem~2 and may use equation (3), with
\begin{equation}
 \overline U(z)=-e^2\int_0^\infty\frac{R_{00}^2(\rho)\rho}{\sqrt{\rho^2+z^2}}\,d\rho.
 \tag{7}
\end{equation}
By using the radial function $R_{00}$ in this expression, we restrict the discussion below to the energy levels of the longitudinal motion that correspond to the lowest Landau level ($\hbar\omega_H/2$) of the transverse motion.

For the ground level, the conditions invoked in the solution of problem 1 in \S~45 are satisfied: since $\chi_0(z)$ has no nodes, it does not vanish at $z=0$. Formula (6), which follows from that solution, is therefore applicable. The ground-state wave function extends over distances $|z|\lesssim a_B$ and varies slowly across this region. The logarithmically divergent integral is consequently cut off at its upper end at distances $|z|\sim a_B$, and at its lower end at distances $|z|\sim a_H$ (there $|z|\sim\rho$, and replacing $\sqrt{\rho^2+z^2}$ by $|z|$ in (7) is not
\SourcePage{556}{559}
permissible). The result is

\begin{equation}
 \varepsilon_0=-\frac{2me^4}{\hbar^2}\ln^2\frac{a_B}{a_H}=-\frac{me^4}{2\hbar^2}\ln^2\frac{\hbar^3H}{m^2c|e|^3}.
 \tag{8}
\end{equation}
This formula is said to possess logarithmic accuracy: both the ratio $a_B/a_H$ itself and its logarithm are assumed large, while the numerical factor in the logarithm's argument remains undetermined.



The excited states of the discrete spectrum are obtained by solving the Schr\"odinger equation (3) with $\overline U(z)\approx-e^2/z$ (which follows from (7) when $z\sim a_B\gg\rho$). Under the substitution $\chi=z\varphi(z)$, however, this equation takes the form


\begin{equation}
 -\frac{\hbar^2}{2m}\frac1{z^2}\frac d{dz}\left(z^2\frac{d\varphi}{dz}\right)-\frac{e^2}{z}\varphi=\varepsilon\varphi,
 \tag{9}
\end{equation}
which has the same form as the equation for the radial wave functions of $s$ states in the three-dimensional Coulomb problem. The required levels are therefore given by (36.10):
\begin{equation}
 \varepsilon_n=-\frac{me^4}{2\hbar^2n^2},
 \tag{10}
\end{equation}
where $n=1,2,3,\ldots$ This expression likewise has only logarithmic accuracy: the next correction would be smaller than the leading term only by the factor $1/\ln(a_B/a_H)$.



Equation (9) determines the wave function only for $z>0$; it can be continued into $z<0$ either as $\chi(-z)=\chi(z)$ or as $\chi(-z)=-\chi(z)$. Accordingly, within the approximation considered, the levels (10) are twofold degenerate. This degeneracy is removed, however, in higher orders in $a_H/a_B$.
